\documentclass[11pt]{scrartcl}
\usepackage[utf8]{inputenc}
\usepackage[sexy]{evanchen} 
\usepackage{graphicx}
\graphicspath{{8 Outsource K Math/Gambar/}}

\title{Vieta Jumping}
\author{Azzam and Akeyla}

\begin{document}

\maketitle

\renewcommand*\contentsname{Daftar Isi}
\tableofcontents

\newpage

\section{Pengantar}
Rangkuman: Ketika extreme principle dan dalil vieta bertemu.

\subsection{IMO 1988 P6: Awal kemunculan di kontes resmi}
\begin{problem}
    [IMO 1988] Misal $a, b$ adalah bilangan asli sehingga $ab + 1$ membagi $a^2 + b^2$. Buktikan bahwa $\dfrac{a^2 + b^2}{ab + 1}$ merupakan bilangan kuadrat
\end{problem}

\begin{proof}[Bukti]
    Misal $k = \dfrac{a^2 + b^2}{ab + 1}$. Asumsikan ada nilai $k$ yang bukan kuadrat. 
    Pilih nilai $A, B$ yang memenuhi $k = \dfrac{A^2 + B^2}{AB + 1}$ dan $A + B$ bernilai minimal. Tanpa mengurangi keumuman, asumsikan bahwa $A \geq B$.
    Fix-kan nilai $B$ dan jadikan $A$ sebagai variabel $x$. Dengan substitusi ke persamaan awal, akan diperoleh persamaan kuadrat \[x^2 - kBx + (B^2 - k) = 0\]
    Tentunya, nilai $A$ merupakan salah satu nilai yang mungkin bagi $x$. Akar kedua dari persamaan kuadrat ini dapat diperoleh dengan dalil vieta: $x_2 = kB - A = \dfrac{B^2 - k}{A}$. Persamaan pertama memastikan bahwa $x_2$ merupakan bilangan bulat dan persamaan kedua memastikan bahwa $x_2 \neq 0$.
    Dapat dicek bahwa $x_2 < A$ (kenapa?). Sehingga, $x_2 + B < A + B$ padahal $k$ juga sama dengan $\dfrac{B^2 + x_2^2}{Bx_2 + 1}$
    Sehingga, dapat disimpulkan bahwa dari awal tidak ada nilai $k$ tidak kuadrat yang memenuhi.
\end{proof}

\subsection{Kesimpulan}
Dapat disimpulkan teknik vieta jumping bekerja sebagai berikut:
\begin{enumerate}
    \item Asumsikan sesuatu yang berlawanan dengan yang ingin dibuktikan
    \item Carilah solusi tertentu yang bernilai paling minimal
    \item Tunjukkan dengan dalil vieta, akan ada solusi lain yang nilainya lebih minimal
    \item Asumsi salah, yang ingin dibuktikan terbukti
\end{enumerate}

\section{Soal-soal}
\begin{problem}
    $a, b$ adalah dua bilangan asli yang memenuhi $ab|a^2 + b^2 + 1$. Buktikan bahwa $a^2 + b^2 + 1 = 3ab$
\end{problem}

\begin{problem}
    Carilah semua pasangan bilangan asli $m, n$ sehingga $\dfrac{m}{n} + \dfrac{n}{m}$ merupakan bilangan bulat
\end{problem}

\begin{problem}
    Buktikan bahwa persamaan $a^2 + b^2 + a + b = 5ab$ tidak punya solusi di bilangan asli
\end{problem}

\begin{problem}
    Hitunglah banyak pasangan bilangan nonnegatif $(m, n)$ yang memenuhi $1 \leq n < m \leq 100, n|m^2 + 1, m|n^2 + 1$.
\end{problem}

\begin{problem}
    [Crux Mathematicorum \#1420] Misal $a, b, c$ adalah bilangan positif dengan $0 < a^2 + b^2 - abc < c$. Buktikan bahwa $a^2 + b^2 - abc$ merupakan bilangan kuadrat
\end{problem}

\begin{problem}
    [IMO 2007 P5] Misal $a, b$ adalah dua bilangan asli dengan $4ab - 1|(4a^2 - 1)^2$. Buktikan bahwa $a = b$. 
\end{problem}

\begin{problem}
    Jika $a, b$ adalah bilangan asli sehingga $\dfrac{a^2 + b^2}{ab - 1} = k$ merupakan bilangan asli, maka buktikan bahwa $k = 5$
\end{problem}

\begin{problem}
    [Turki TST 1994] Misal $a, b$ adalah bilangan asli dengan sehingga $ab|a^2 + b^2 + 3$. Buktikan bahwa $\dfrac{a^2 + b^2 + 3}{ab}$ hanya mungkin bernilai $4$ atau $5$
\end{problem}

\begin{problem}
    [AMM 1994] Carilah semua bilangan asli $k$ sedemikian hingga terdapat bilangan-bilangan asli $x, y, z$ yang memenuhi $(x + y + z)^2 = kxyz$
\end{problem}

\begin{problem}
    Buktikan bahwa apabila $P = \dfrac{x^2 + 1}{y^2} + 4$ merupakan bilangan kuadrat dengan $x, y$ bilangan asli, maka haruslah $P = 9$
\end{problem}

\newpage
\section{Solusi}
\begin{enumerate}
    \item Tulis $a^2 + b^2 + 1 = kab$ untuk suatu bilangan asli $k$. Asumsikan ada nilai $a$ dan $b$ sehingga $k \neq 3$. Pilih pasangan $(a, b)$ yang jumlahnya paling kecil. WLOG, $a \geq b$. Tetapkan nilai $b$ dan jadikan $a$ sebagai variabel $x$. Substitusikan ke persamaan awal dan didapat persamaan kuadrat \[x^2 - kbx + b^2 + 1 = 0.\] Persamaan kuadrat tersebut memiliki akar yang salah satunya sama dengan $a$. Akar keduanya adalah $x_2 = kb - a = \frac{b^2 + 1}{a}$. Persamaan pertama memastikan bahwa $x_2$ adalah bilangan bulat dan persamaan kedua memastikan bahwa $x_2$ positif. 
    
    Amatilah bahwa $x_2 = \frac{b^2 + 1}{a} \leq \frac{a^2 + 1}{a} = a + \frac{1}{a}$. Selain dari kasus $a = b = 1$, kita punya $x_2 \leq a$. Jika $a = b = 1$ maka tentunya $k = 3$. Kemudian, jika $x_2 = a$ maka $\frac{b^2 + 1}{a} = x_2 = a$ yang membuat $b^2 + 1 = a^2$. Persamaan ini tidak memiliki solusi. Sehingga, dapat disimpulkan bahwa $x_2 < a$ dan $(x_2, b)$ memenuhi persamaan awal. Hal ini kontradiktif dengan asumsi bahwa $(a, b)$ adalah pasangan solusi dengan jumlah yang paling kecil. Jadi, dapat disimpulkan bahwa $k$ harus sama dengan $3$.
    \item Tentunya, $(r, r)$ merupakan solusi dengan $\frac{m}{n} + \frac{n}{m} = 2$. Dengan AM-GM, didapat bahwa nilai tersebut minimal sama dengan $2$ dan kesamaan hanya dapat tercapai saat $\frac{m}{n} = \frac{n}{m}$, yakni $m = n$. 

    Asumsikan ada solusi $(m, n)$ dengan $\frac{m}{n} + \frac{n}{m} = k > 2$. WLOG, $m > n$. Tinjau solusi dengan jumlah terkecil. Tetapkan nilai $n$ dan variasikan nilai $m$. Dengan menyusun ulang persamaan, didapat persamaan kuadrat \[x^2 - kxn + n^2 = 0.\] Misal $x_2$ adalah akar salah satunya lagi. Amati bahwa $x_2 = kx - m = \frac{n^2}{m}$. Tentunya, $x_2 =\frac{n^2}{m} < \frac{m^2}{m} = m$ yang membuat $(x_2, n)$ adalah solusi dengan jumlah yang lebih kecil. Sehingga, terjadi kontradiksi dan $k$ harus sama dengan $2$.
    \item Asumsikan ada solusi dengan jumlah $a + b$ yang paling kecil. WLOG, $a > b$ (kasus $a = b$ tentunya tidak memenuhi). Tetapkan nilai $b$ dan variasikan nilai $a$. Persamaan kuadrat \[x^2 + (1 - 5b)x + b^2 + b\] memiliki akar $a$ dan $x_2 = 5b - 1 - a = \frac{b^2 + b}{a}$. Tentunya, $x_2 = \frac{b^2 + b}{a} < \frac{a^2 + a}{a} = a + 1$. Sehingga, $x_2 \geq a$. Namun, dikarenakan $x_2 = a = b$ tidak memenuhi persamaan, dapat disimpulkan bahwa $x_2 < a$ dan kita memperoleh solusi yang lebih kecil, kontradiksi. Jadi, sama sekali tidak ada solusi.
    \item Tentunya, $m$ dan $n$ harus relatif prima dan $mn|(m^2 + 1)(n^2 + 1) = m^2n^2 + m^2 + n^2 +1$. Akibatnya, $mn$ membagi $m^2 + n^2 + 1$ dan dengan soal pertama, didapat $m^2 + n^2 + 1 = 3mn$. Tentunya, tanpa memperhatikan syarat dari $m > n$, solusi dengan $m + n$ terkecil adalah $(1, 1)$. Dengan melihat solusi pada soal pertama, apabila $(m, n)$ dengan $m > n$ merupakan solusi maka $(3n - m, n)$ juga merupakan solusi. Jadi, semua solusi dapat dihasilkan dengan barisan rekursif $(m_k, n_k) = (3n_k - m_k, n_k)$. Dimulai dari $(1, 1)$, bisa didapat secara berturut-turut $(2, 1), (5, 2), (13, 5), (34, 13), (89, 34)$. Sehingga, akan didapat $10$ solusi.
    \item Asumsikan ada nilai $k = a^2 + b^2 - abc$ yang bukan merupakan bilangan kuadrat. Tentunya apabila $a = b$ maka $0 < (2 - c)a^2 < c$ yang hanya memiliki solusi $c = 1$ dan $a^2 + b^2 - abc = a^2$ merupakan bilangan kuadrat. Jadi, WLOG $a > b$ dan ambil solusi dengan jumlah $a + b$ minimal. Tetapkan nilai $b$ dan variasikan nilai $a$. Persamaan kuadrat \[x^2 - bcx + b^2 -k = 0\] memiliki solusi $a$ dan $x_2 = bc - a = \frac{b^2 - k}{a}$. Tentunya, $x_2 = \frac{b^2 - k}{a} < \frac{b^2}{a} < \frac{a^2}{a} = a$ yang membuat $(x_2, b)$ adalah solusi dengan jumlah yang lebih kecil. Dapat disimpulkan bahwa tidak ada solusi $(a, b)$ dengan $k$ bukan merupakan bilangan kuadrat.
    \item Amati bahwa $4ab - 1$ juga akan membagi $b^2(4a^2 - 1)^2 - (4ab - 1)(4a^3b - 2ab + a^2) = (a - b)^2$. Asumsikan ada solusi dengan $a \neq b$. Tulis $k = \dfrac{(a - b)^2}{4ab - 1} > 0$. WLOG, $a > b$ dan pilih salah satu solusi dengan jumlah $a + b$ paling kecil. Tetapkan nilai $b$ dan variasikan nilai $a$. Didapat persamaan kuadrat \[x^2 - (2b + 4kb)x + b^2 + k = 0\] yang memiliki solusi $a$ dan $x_2 = 2b + 4kb - a = \frac{b^2 + k}{a}$. Dengan asumsi bahwa $a + b$ memiliki nilai terkecil yang mungkin, didapat $\frac{b^2 + k}{a} > a$ yang ekivalen dengan $k > a^2  - b^2$. Jadi, $\frac{(a - b)^2}{4ab - 1} = k > a^2 - b^2$. Sehingga, $a - b > (a + b)(4ab - 1)$. Namun, ketaksamaan terakhir ini tidak memiliki solusi. Dapat disimpulkan bahwa memang tidak ada solusi dengan $a = b$.
    \item Tentunya, jika $a = b$ maka $k = \frac{2a^2}{a^2 - 1} = 2 + \frac{2}{a^2 - 1}$ yang tidak pernah merupakan bilangan asli. Asumsikan ada solusi dengan $k \neq 5$ dan WLOG $a > b$. Pilih solusi dengan jumlah $a + b$ terkecil. Tetapkan nilai $b$ dan variasikan nilai $a$. Persamaan kuadrat \[x^2 - kbx + b^2 + k = 0\] memiliki akar $a$ dan $x_2 = kb - a = \frac{b^2 + k}{a}$. Dengan asumsi bahwa $x_2 \geq a$ (karena $(a, b)$ adalah solusi dengan jumlah terkecil), maka $\frac{b^2 + k}{a} > a$ yang membuat $k > a^2 - b^2 \geq a^2 - (a - 1)^2 = 2a - 1$. Jadi, dapat disimpulkan bahwa $a^2 + b^2 = k(ab - 1) > (2a - 1)(ab - 1)$. Lebih lanjut, $a^2 > a^2 + b^2 > (2a - 1)(ab - 1) > a\cdot a(b - 1)$ yang memberikan solusi $b = 1$. Dapat dicek bahwa $b = 1$ memberikan $k = \frac{a^2 + 1}{a - 1} = a + 1 + \frac{2}{a - 1}$ yang hanya memiliki solusi $a = 2$ atau $3$ dan $k$ harus sama dengan $5$. Jadi, $k$ harus sama dengan $5$.
    \item Asumsikan $k \neq 4$ atau $5$. Ambil solusi dengan jumlah $a + b$ minimal. WLOG, $a > b$. Amati bahwa $a = b$ hanya dapat tercapai saat $a = b = 1$ dan $k = 5$. Tentunya, polinomial \[x^2 - kbx + b^2 + 3 = 0\] memiliki akar $a$ dan $x_2 = kb - a = \frac{b^2 + 3}{a}$. Dikarenakan $(a, b)$ adalah solusi minimal maka $a \leq x_2 = \frac{b^2 + 3}{a}$ yang membuat $3 \geq a^2 - b^2 \geq a^2 - (a - 1)^2 = 2a - 1$. Artinya, $a = 1$ atau $2$. Namun, karena $a > b$ maka kemungkinan yang tersisa adalah $a = 2$ dan $b = 1$ yang memberikan solusi $k = 5$. Jadi, $k = 4$ atau $5$.
    \item Untuk sebuah bilangan asli $k$, pilih solusi $(X, Y, Z)$ dengan jumlah yang minimal. WLOG, $X \geq Y \geq Z$ dan variasikan nilai $X$ dengan menjaga nilai $Y$ dan $Z$ tetap. Polinomial \[x^2 + (2(Y + Z) - kYZ)x + (Y + Z)^2 = 0\] memiliki akar $X$ dan $x_2 = kYZ - 2(Y + Z) = \frac{(Y + Z)^2}{X}$. Dikarenakan $X + Y + Z$ minimal, haruslah $\frac{(Y + Z)^2}{X} = x_2 \leq X$  yang memberikan $Y + Z \leq X$. Substitusikan ketaksamaan ini ke persamaan awal dan didapat 
    \begin{align*}
        k &= \dfrac{(X + Y + Z)^2}{XYZ}\\
        &= \dfrac{(X+Y + Z)}{XZ} \cdot \dfrac{(X + Y + Z)}{Y}\\
        &\leq \dfrac{(X+ X + X)}{XZ} \cdot \dfrac{(Y + Z) + Y + Z}{Y}\\
        &\leq \dfrac{3}{Z} \cdot \dfrac{4Y}{Y}\\
        &= \dfrac{12}{Z}
    \end{align*}
    Dikarenakan $Z$ adalah bilangan asli maka $k \leq 12$. Jika $Z = 1$ maka $X \leq Y + Z = Y + 1$ membuat $X = Y$ atau $X = Y + 1$. Substitusikan nilai ini ke persamaan awal dan akan didapat nilai $k = 5, 6, 8$ dan $9$. 

    Jika $Z > 1$ maka $k < 6$. Hanya perlu dicek apakah $k = 1, 2, 3, 4$ memenuhi. Dapat dicek manual bahwa $(9, 9, 9), (8, 4, 4), (3, 3, 3), (4, 2, 2)$ secara berturut-turut memberikan solusi $k = 1, 2, 3, 4$. Jadi, $k = 1, 2, 3, 4, 5, 6, 8, 9$.
    \item Tulis $k^2 = \frac{x^2 + 1}{y^2} + 4$. Manipulasi aljabar akan menghasilkan $x^2 - (k^2 - 4)y^2 = -1$. Kalikan kedua ruas dengan $4$ dan didapat $u^2 - (k^2 - 4)v^2 = -4$ dengan $u = 2x$ dan $v = 2y$. Tulis $m = \frac{u + kv}{2}$. Didapat $m^2 + v^2 + 1 = kmv$. Ini adalah soal pertama dan dapat diselesaikan dengan Vieta Jumping yang menghasilkan $k = 3$ dan $P = k^2 = 9$.
\end{enumerate}

\end{document}